%! Graphing Exponential Functions and Transformations — Unit 7, Lesson 7.2 — Group Activity (Answer Key)
\documentclass[10pt]{article}
\usepackage{algebra2-article}
\usepackage{algebra2-key}   % pulls in -boxes; do NOT also load -boxes
\newcommand{\TallMath}[1]{$\displaystyle #1\rule[-1.4em]{0pt}{3.2em}$}
\newcommand{\lgrid}[4]{%
  \draw[step=1, linegray!40, very thin] (#1,#3) grid (#2,#4);
  \draw[->, semithick, charcoal] (#1,0)--(#2,0) node[below right, font=\scriptsize]{$x$};
  \draw[->, semithick, charcoal] (0,#3)--(0,#4) node[left, font=\scriptsize]{$y$};}
% g(x) = a*b^(x-h) + k plotted as a*exp((ln b)(x-h))+k.
% #1=a  #2=ln(b)  #3=h  #4=k  #5=xmin  #6=xmax
% ln2 = 0.693147   ln3 = 1.098612   ln(1/2) = -0.693147
\newcommand{\expgraph}[6]{\draw[very thick, forest, domain=#5:#6, samples=140, smooth]
  plot (\x,{(#1)*exp((#2)*((\x)-(#3)))+(#4)});}
\newcommand{\hasym}[3]{\draw[navy, dashed, semithick] (#1,#3)--(#2,#3);}

\begin{document}

\pageheader{Unit 7, Lesson 7.2}{Group Activity --- Answer Key}
\namepartnerperiod

\vspace{0.08in}

\begin{headlinebox}{forestbg}
\small\textbf{Move It --- and Say What Moved.} Every problem below starts from a parent $y=b^{x}$ and
does something to it. For each one, your partner should be able to stop you and ask the two questions
that decide everything: \textbf{did the asymptote move?} and \textbf{what is the range now?} If a
transformation did not move the dashed line, it did not change the range either.
\end{headlinebox}

\vspace{0.06in}

\begin{tcolorbox}[colback=white, colframe=black!40, title=\textbf{Tier R --- Remediate},
    fonttitle=\bfseries, arc=1.5mm, left=3mm, right=3mm, top=2mm, bottom=2mm, breakable]
\textbf{Part 1 --- What moved?} Every row starts from the parent $y=3^{x}$, whose asymptote is $y=0$
and whose range is $(0,\infty)$.
\begin{center}
\renewcommand{\arraystretch}{1.6}
\begin{tabularx}{\linewidth}{@{}l X l X@{}}
\hline
\textbf{Function} & \textbf{Move(s) from the parent} & \textbf{Asymptote} & \textbf{Range} \\
\hline
$y=3^{x}+4$ & \ans{up $4$} & \ans{$y=4$} & \ans{$(4,\infty)$} \\
$y=3^{\,x-2}$ & \ans{right $2$} & \ans{$y=0$} & \ans{$(0,\infty)$} \\
$y=3^{x}-5$ & \ans{down $5$} & \ans{$y=-5$} & \ans{$(-5,\infty)$} \\
$y=-3^{x}$ & \ans{flip over the $x$-axis} & \ans{$y=0$} & \ans{$(-\infty,0)$} \\
\hline
\end{tabularx}
\end{center}
Two of those four rows have the same asymptote as the parent. Which two, and what do they have in
common? \ans{Rows 2 and 4 --- neither one is a \emph{vertical shift}.}

\vspace{6pt}
\textbf{Part 2 --- One function, all the way through.} Let $g(x)=2^{x}+3$.
\begin{enumerate}[label=(\alph*), leftmargin=1.9em, itemsep=6pt]
  \item Parent function: $y=$ \ans{$2^{x}$} \quad Move: \ans{up $3$}
  \item Horizontal asymptote: $y=$ \ans{3}
  \item Fill in the table. (Remember $2^{-1}=\frac12$ and $2^{-2}=\frac14$.)
    \begin{center}
    \renewcommand{\arraystretch}{1.35}
    \begin{tabular}{|c|c|c|c|c|c|}
    \hline
    $x$ & $-2$ & $-1$ & $0$ & $1$ & $2$ \\ \hline
    $g(x)$ & \ans{3.25} & \ans{3.5} & \ans{4} & \ans{5} & \ans{7} \\ \hline
    \end{tabular}
    \end{center}
  \item $y$-intercept: \ans{$(0,4)$} \quad Range: \ans{$(3,\infty)$}
  \item Every output in your table is bigger than $3$. Will \emph{any} input ever give an output of
        exactly $3$? Why not? \ans{No --- $2^{x}$ is always positive, never $0$, so $2^{x}+3>3$ always.}
  \item Does this graph have an $x$-intercept? Explain using the asymptote.
        \ans{No. The whole graph sits above $y=3$, which is already above the $x$-axis.}
\end{enumerate}
\end{tcolorbox}

\begin{tcolorbox}[colback=white, colframe=black!40, title=\textbf{Tier A --- Approaching Proficiency},
    fonttitle=\bfseries, arc=1.5mm, left=3mm, right=3mm, top=2mm, bottom=2mm, breakable]
\textbf{Part 1 --- Match, then read.} Each graph below is a transformation of $y=2^{x}$. The dashed
line is the horizontal asymptote.
\begin{center}
\begin{tikzpicture}[scale=0.28]
  % P: y = 2^x + 2
  \begin{scope}
    \lgrid{-2}{5}{-4}{8}
    \begin{scope}\clip (-2,-4) rectangle (5,8); \expgraph{1}{0.693147}{0}{2}{-2}{5} \end{scope}
    \hasym{-2}{5}{2}
    \node[charcoal, font=\small] at (1.5,-5.6) {\textbf{P}};
  \end{scope}
  % Q: y = 2^x - 2
  \begin{scope}[xshift=9.5cm]
    \lgrid{-2}{5}{-4}{8}
    \begin{scope}\clip (-2,-4) rectangle (5,8); \expgraph{1}{0.693147}{0}{-2}{-2}{5} \end{scope}
    \hasym{-2}{5}{-2}
    \node[charcoal, font=\small] at (1.5,-5.6) {\textbf{Q}};
  \end{scope}
  % R: y = 2^(x-2)
  \begin{scope}[xshift=19cm]
    \lgrid{-2}{5}{-4}{8}
    \begin{scope}\clip (-2,-4) rectangle (5,8); \expgraph{1}{0.693147}{2}{0}{-2}{5} \end{scope}
    \hasym{-2}{5}{0}
    \node[charcoal, font=\small] at (1.5,-5.6) {\textbf{R}};
  \end{scope}
  % S: y = 2^(-x)
  \begin{scope}[xshift=28.5cm]
    \lgrid{-2}{5}{-4}{8}
    \begin{scope}\clip (-2,-4) rectangle (5,8); \expgraph{1}{-0.693147}{0}{0}{-2}{5} \end{scope}
    \hasym{-2}{5}{0}
    \node[charcoal, font=\small] at (1.5,-5.6) {\textbf{S}};
  \end{scope}
\end{tikzpicture}
\end{center}
\vspace{2pt}
Write the letter of the matching graph, then fill in the last two columns from the \emph{equation}.
\begin{center}
\renewcommand{\arraystretch}{1.55}
\begin{tabularx}{\linewidth}{@{}l c X X@{}}
\hline
\textbf{Equation} & \textbf{Graph} & \textbf{Asymptote} & \textbf{Range} \\
\hline
$y=2^{x}+2$ & \ans{P} & \ans{$y=2$} & \ans{$(2,\infty)$} \\
$y=2^{x}-2$ & \ans{Q} & \ans{$y=-2$} & \ans{$(-2,\infty)$} \\
$y=2^{\,x-2}$ & \ans{R} & \ans{$y=0$} & \ans{$(0,\infty)$} \\
$y=2^{-x}$ & \ans{S} & \ans{$y=0$} & \ans{$(0,\infty)$} \\
\hline
\end{tabularx}
\end{center}
Exactly one of those four graphs has an $x$-intercept. Which one, and why is it the only one?
\ans{Q --- it is the only one whose asymptote sits \emph{below} the $x$-axis, so its range includes $0$.}

\vspace{6pt}
\textbf{Part 2 --- Backwards: from the graph to the equation.} The curve has the form
$y=a\,b^{x}+k$.
\begin{center}
\begin{minipage}[c]{0.42\linewidth}
\centering
\begin{tikzpicture}[scale=0.40]
  \lgrid{-3}{4}{-3}{11}
  \begin{scope}\clip (-3,-3) rectangle (4,11); \expgraph{3}{0.693147}{0}{-1}{-3}{4} \end{scope}
  \hasym{-3}{4}{-1}
  \node[navy, font=\scriptsize, left] at (-1.2,-1.7) {$y=-1$};
  \fill[charcoal] (0,2) circle (3pt) node[left, font=\scriptsize]{$(0,2)$};
  \fill[charcoal] (1,5) circle (3pt) node[right, font=\scriptsize]{$(1,5)$};
\end{tikzpicture}
\end{minipage}
\begin{minipage}[c]{0.54\linewidth}
Step 1 --- the dashed line gives $k=$ \ans{$-1$} \\[6pt]
Step 2 --- the $y$-intercept is $a+k$, so $a=$ \ans{3} \\[6pt]
Step 3 --- use $(1,5)$: \; $5=$ \ans{3}$\cdot b+$ \ans{$(-1)$}, so $b=$ \ans{2} \\[6pt]
Equation: $y=$ \ans{$3(2)^{x}-1$} \\[6pt]
Range: \ans{$(-1,\infty)$}
\end{minipage}
\end{center}

\vspace{4pt}
\textbf{Justify.} Nadia says $y=-5^{x}$ and $y=5^{-x}$ are ``the same thing, just written
differently --- they both have a negative.'' She is wrong. For \emph{each} function, say which axis
it reflects over, give its range, and say whether it is growth or decay. Then state in one sentence
what Nadia is actually confusing.
\ansline{$y=-5^{x}$ reflects over the \textbf{$x$}-axis. Range $(-\infty,0)$. The base is still $5$,
so it is still a growth function --- the graph falls only because $a<0$.\\[3pt]
$y=5^{-x}$ reflects over the \textbf{$y$}-axis. Range $(0,\infty)$, unchanged. It equals
$\left(\frac15\right)^{x}$, so it is \textbf{decay}.\\[3pt]
Nadia is confusing a negative on the \emph{output} (which flips the graph over the $x$-axis and moves
the range to the other side of the asymptote) with a negative on the \emph{input} (which flips it over
the $y$-axis and changes the base).}
\end{tcolorbox}

\begin{tcolorbox}[colback=white, colframe=black!40, title=\textbf{Tier E --- Extension},
    fonttitle=\bfseries, arc=1.5mm, left=3mm, right=3mm, top=2mm, bottom=2mm, breakable]
\textbf{Part 1 --- A horizontal shift in disguise.} Exponentials do something no other family in this
course does. Use the exponent rules from Unit 6.
\begin{enumerate}[label=(\alph*), leftmargin=1.9em, itemsep=6pt]
  \item Rewrite $2^{\,x-3}$ as a product: \; $2^{\,x-3}=2^{x}\cdot 2^{\,\square}$ where $\square=$
        \ans{$-3$}, \; and $2^{\,\square}=$ \ans{$\frac18$}. \\[3pt]
        So $2^{\,x-3}=$ \ans{$\frac18$} $\cdot\, 2^{x}$.
  \item Check it in a table.
    \begin{center}
    \renewcommand{\arraystretch}{1.35}
    \begin{tabular}{|c|c|c|c|c|}
    \hline
    $x$ & $3$ & $4$ & $5$ & $6$ \\ \hline
    $2^{\,x-3}$ & \ans{1} & \ans{2} & \ans{4} & \ans{8} \\ \hline
    $\frac18\cdot 2^{x}$ & \ans{1} & \ans{2} & \ans{4} & \ans{8} \\ \hline
    \end{tabular}
    \end{center}
  \item So for an exponential, shifting \emph{right} $3$ produces the same graph as a vertical
        \ans{compression} by a factor of $\frac18$. Does either move change the asymptote?
        \ans{No --- still $y=0$.}
  \item Try the same trick on the parabola: is $(x-3)^{2}$ equal to some fixed number times $x^{2}$?
        Test it at $x=4$ and $x=5$ and explain what goes wrong.
        \ansline{At $x=4$: $(4-3)^{2}=1$ while $x^{2}=16$, a ratio of $\frac{1}{16}$. At $x=5$:
        $4$ versus $25$, a ratio of $\frac{4}{25}$. The ratio is not fixed, so no.\\[3pt]
        The trick works only for exponentials because there the input sits in the \emph{exponent},
        where subtracting $3$ factors out as the constant $2^{-3}$.}
\end{enumerate}

\vspace{5pt}
\textbf{Part 2 --- All four moves at once.} Let $g(x)=-2(3)^{\,x-1}+5$.
\begin{center}
\renewcommand{\arraystretch}{1.5}
\begin{tabularx}{\linewidth}{@{}l X@{}}
\hline
Parent function & \ans{$y=3^{x}$} \\
The moves, in words & \ans{right $1$; stretch vertically by $2$ and flip over the $x$-axis; up $5$} \\
Horizontal asymptote & \ans{$y=5$} \\
Range & \ans{$(-\infty,5)$} \\
$y$-intercept ($g(0)$, exact) & \ans{$\left(0,\tfrac{13}{3}\right)$, since $-2\cdot\tfrac13+5=\tfrac{13}{3}$} \\
\hline
\end{tabularx}
\end{center}
Does $g$ have an $x$-intercept? Answer using the signs of $a$ and $k$ --- not by graphing.
\ans{Yes. $a=-2<0$ and $k=5>0$ are opposite signs, so the range $(-\infty,5)$ contains $0$.}

\vspace{5pt}
\textbf{Part 3 --- When the asymptote is a real thing.} A cup of coffee is poured at
$200^{\circ}$F into a room held at $68^{\circ}$F. Its temperature after $t$ minutes is
\[
  T(t) = 68 + 132(0.85)^{t}.
\]
\begin{center}
\renewcommand{\arraystretch}{1.3}
\begin{tabular}{|c|c|c|c|c|c|c|}
\hline
$t$ (min) & $0$ & $4$ & $8$ & $12$ & $16$ & $20$ \\ \hline
$T(t)$ ($^{\circ}$F) & $200$ & $136.9$ & $104.0$ & \ans{86.8} & $77.8$ & $73.1$ \\ \hline
\end{tabular}
\end{center}
\begin{enumerate}[label=(\alph*), leftmargin=1.9em, itemsep=6pt]
  \item Identify $a$, $b$, and $k$: \; $a=$ \ans{132} \quad $b=$ \ans{0.85} \quad
        $k=$ \ans{68}
  \item Fill the missing table entry by evaluating $T(12)$. \ans{$68+132(0.1422)\approx 86.8^{\circ}$F}
  \item The horizontal asymptote is $y=$ \ans{68}. What does that number \emph{mean} about
        this coffee? \ans{Room temperature --- the coffee cools toward $68^{\circ}$ but never reaches it.}
  \item Algebraically the range of $T$ is $(68,\infty)$. But the temperatures this coffee actually
        takes are only $(68,200]$. Explain the difference.
        \ansline{The algebraic range allows every real $t$, including negative ones. The story starts
        the clock at $t=0$, and the coffee is hottest at that instant, so $200$ is the largest output
        the situation ever produces --- and it \emph{is} reached, which is why that end is a bracket.}
  \item At what time is the coffee exactly $68^{\circ}$F? \ans{Never --- $68$ is the asymptote.}
  \item Between which two table times does the coffee pass $100^{\circ}$F? \ans{between $8$ and $12$ min} \\[3pt]
        To find that moment \emph{exactly} you would have to solve $132(0.85)^{t}=32$. Try rewriting
        both sides with a common base --- what goes wrong?
        \ansline{Dividing gives $(0.85)^{t}=\frac{32}{132}\approx 0.2424$, and there is no base that
        writes both $0.85$ and $0.2424$ as clean powers. The common-base method cannot reach this one ---
        we need a new tool.}
\end{enumerate}
\end{tcolorbox}

\begin{teachernote}
\textbf{Tier R Part 1's closing question is the whole activity.} Rows 2 and 4 keep the parent's
asymptote because neither is a \emph{vertical shift}; rows 1 and 3 move it because both are. Students
who answer that correctly can be sent straight to Tier A.

\textbf{Tier A's matching is diagnostic, not busywork.} Q is the only graph with an $x$-intercept, and
the reason --- its asymptote is below the $x$-axis --- is the payoff of Lesson 7.0's ``no zeros ever.''
Watch for pairs who match P and Q by ``it looks higher'' instead of reading the dashed line; make them
name the asymptote out loud before writing a letter. In Part 2, the signature error is $a=2$ (grabbing
the $y$-intercept) instead of $a=3$; the $y$-intercept is $a+k$, and here $k$ is \emph{negative}, so
$a$ comes out \emph{larger} than the intercept. That surprise is worth a whole-class pause.

\textbf{The Justify item is the assessed idea of the lesson} and is worth more time than its length
suggests. Both students' functions contain a minus sign, but one negates the \emph{output} and the
other negates the \emph{input}. Require all three pieces for each function --- axis, range, family ---
because a student can get ``$x$-axis'' right and still say $-5^{x}$ is decay. It is not: the base is
$5$, and Lesson 7.0's rule is that growth versus decay is read off $b$ against $1$, never off $a$.

\textbf{Tier E Part 1 is genuinely new mathematics, not review.} $2^{\,x-3}=\frac18\cdot 2^{x}$ says a
horizontal shift and a vertical compression are the \emph{same transformation} for this family --- the
exponential analogue of the input/output equivalence students met for radicals in Lesson 6.4. Part (d)
is what keeps it honest: the trick fails for $y=x^{2}$ because there the input is not in an exponent.
Do not let a group skip (d).

\textbf{Tier E Part 3 is where $k$ stops being decoration.} The asymptote is room temperature, a
physical fact, and part (d) is a contextual-domain conversation worth having: an algebraic range and a
realistic range are different objects. Part (f) hits Lesson 7.4's wall on purpose --- $0.85$ and
$\frac{8}{33}$ share no convenient base. Leave it standing; do not mention logarithms. Lesson 7.3 will
solve the ones that \emph{do} cooperate, and Unit 8 will come back for this one.

\textbf{Pacing.} Tier E Part 3 is the piece to protect if time runs short; Part 1 can go home as an
extension.
\end{teachernote}

\end{document}
